\begin{table}[t]
\centering
\small
\caption{Repeated-runs comparison: each program run 30 times per data set with different seeds (8 chains of 15,000 iterations). Per parameter, the 30 posterior means and SDs from jmbayes are compared with the 30 from JMbayes2. Shown equivalent: two one-sided Welch tests at 5\%, for the posterior mean within $\pm 0.1$ and $\pm 0.05$ posterior SD, and for the posterior SD within $\pm 10\%$. Shown to differ: Welch's $t$-test (mean, SD) and the two-sample Kolmogorov--Smirnov test (whole run distribution of the posterior mean), after Benjamini--Hochberg correction at 5\% false discovery rate. 80\% power: parameters for which 30 runs per program give at least 80\% power to show equivalence of the mean within $\pm 0.1$ SD.}\label{tab:multi}
\vspace{4pt}
\begin{tabular}{lrrrrrrrr}
\toprule & & \multicolumn{3}{c}{Shown equivalent} & \multicolumn{3}{c}{Shown to differ} & \\
\cmidrule(lr){3-5}\cmidrule(lr){6-8}
& & \multicolumn{2}{c}{Mean} & SD & & & Run & 80\% \\
Data set & Param. & $\pm 0.1$ & $\pm 0.05$ & $\pm 10\%$ & Mean & SD & distr. & power \\ \midrule
PBC (312 patients, 3 markers) & 50 & 50 & 43 & 50 & 0 & 0 & 0 & 50 \\
AIDS (467 patients) & 22 & 22 & 22 & 22 & 0 & 0 & 0 & 22 \\
Liver cirrhosis (488 patients) & 22 & 22 & 22 & 22 & 0 & 0 & 0 & 22 \\
\bottomrule
\end{tabular}
\end{table}
